\documentclass[11pt]{article}
\usepackage[utf8]{inputenc}
\usepackage{amsmath,amssymb,booktabs,geometry,verbatim}
\usepackage{CJKutf8}
\geometry{margin=1in}

\title{CSP-J 2021 (Junior)}
\date{2021-10-23, 8:30--12:00 (3.5 hours)}

\begin{document}
\maketitle

\noindent
Translated from the official Chinese statements of
the second round of CSP-J 2021 (CCF non-professional software capability
certification, junior level).

\section*{Problem 1: Distributing Candy
  (\begin{CJK}{UTF8}{gbsn}分糖果\end{CJK}, candy)}

Time limit: 1 s.
Memory limit: 512 MB.
Total points: 100.

\subsection*{Background}

The children of Red Sun Kindergarten are starting to distribute candy!

\subsection*{Description}

Red Sun Kindergarten has $n$ children, and you are one of them.
It is guaranteed that $n \ge 2$.

One day you discover infinitely many pieces of candy in the back garden of
the kindergarten, and you plan to take some candy back to distribute among
the children of the kindergarten.

Since you are just an ordinary kindergarten child, your strength is
limited, so you can take back at most $R$ pieces of candy.

But taking too few would not be enough to distribute, so you must take back
at least $L$ pieces of candy.
It is guaranteed that $n \le L \le R$.

In other words, if you take $k$ pieces of candy, then you must guarantee
that $L \le k \le R$.

If you take $k$ pieces of candy, you will put these $k$ pieces into a
basket and ask everyone to distribute the candy according to the following
scheme:
as long as the basket contains \textbf{no fewer than} $n$ pieces of candy,
all $n$ children of the kindergarten (including yourself) each take
\textbf{exactly} one piece of candy from the basket,
until the number of pieces in the basket is \textbf{fewer than} $n$.
At this moment all the candy remaining in the basket belongs to you ---
these pieces are \textbf{your reward for carrying the candy}.

As a high-quality kindergarten child, you want the number of pieces that
are \textbf{your reward for carrying the candy}
(\textbf{not the total number of pieces you end up with!})
to be as large as possible;
therefore you need to write a program that reads $n, L, R$ in order and
outputs the maximum number of pieces of candy you can obtain as
\textbf{your reward for carrying the candy}.

\subsection*{Input}

Read the input from the file \texttt{candy.in}.

The input consists of one line containing three positive integers
$n, L, R$,
denoting the number of children, the lower bound on the number of pieces
of candy, and the upper bound, respectively.

\subsection*{Output}

Write the output to the file \texttt{candy.out}.

Output one line containing one integer:
the maximum number of pieces of candy you can obtain as
\textbf{your reward for carrying the candy}.

\subsection*{Sample 1 input}

\begin{verbatim}
7 16 23
\end{verbatim}

\subsection*{Sample 1 output}

\begin{verbatim}
6
\end{verbatim}

\subsection*{Explanation of sample 1}

Take $k = 20$ pieces of candy and put them into the basket.

The basket now contains $20 \ge n = 7$ pieces, so every child gets one
piece;

the basket now contains $13 \ge n = 7$ pieces, so every child gets one
piece;

the basket now contains $6 < n = 7$ pieces, so these 6 pieces are
\textbf{your reward for carrying the candy}.

It is easy to see that the number of pieces you obtain as
\textbf{your reward for carrying the candy} cannot exceed 6
(otherwise the number of pieces left in the basket would still be no fewer
than $n$, and every child would have to take another piece),
so the answer is 6.

\subsection*{Sample 2 input}

\begin{verbatim}
10 14 18
\end{verbatim}

\subsection*{Sample 2 output}

\begin{verbatim}
8
\end{verbatim}

\subsection*{Explanation of sample 2}

It is easy to see that when the number $k$ of pieces you take satisfies
$14 = L \le k \le R = 18$,
after every child gets one piece the remaining $k - 10$ pieces are always
the number of pieces that are
\textbf{your reward for carrying the candy},
so taking $k = 18$ pieces is optimal, and the answer is 8.

\subsection*{Sample 3}

See \texttt{candy/candy3.in} and \texttt{candy/candy3.ans} in the
contestant's directory.

\subsection*{Constraints}

\begin{center}
\begin{tabular}{cccc}
\toprule
Test point & $n \le$ & $R \le$ & $R - L \le$ \\
\midrule
1  & 2      & 5      & 5      \\
\midrule
2  & 5      & 10     & 10     \\
\midrule
3  & $10^3$ & $10^3$ & $10^3$ \\
\midrule
4  & $10^5$ & $10^5$ & $10^5$ \\
\midrule
5  & $10^3$ & $10^9$ & 0      \\
\cmidrule(lr){1-1}\cmidrule(lr){4-4}
6  &        &        & $10^3$ \\
\cmidrule(lr){1-2}\cmidrule(lr){4-4}
7  & $10^5$ &        & $10^5$ \\
\cmidrule(lr){1-2}\cmidrule(lr){4-4}
8  & $10^9$ &        & $10^9$ \\
\cmidrule(lr){1-1}
9  &        &        &        \\
\cmidrule(lr){1-1}
10 &        &        &        \\
\bottomrule
\end{tabular}
\end{center}

In the original table the cell $n \le 10^3$ spans test points 5--6,
the cell $R \le 10^9$ spans test points 5--10,
and the cells $n \le 10^9$ and $R - L \le 10^9$ span test points 8--10.

For all data, it is guaranteed that $2 \le n \le L \le R \le 10^9$.

\section*{Problem 2: Insertion Sort
  (\begin{CJK}{UTF8}{gbsn}插入排序\end{CJK}, sort)}

Time limit: 1 s.
Memory limit: 512 MB.
Total points: 100.

\subsection*{Description}

Insertion sort is a very common and simple sorting algorithm.
Little Z is a freshman, and today Teacher H has just taught the insertion
sort algorithm in class.

Suppose comparing two elements takes $O(1)$ time;
then insertion sort can sort an array of length $n$ in $O(n^2)$ time.
Suppose the $n$ numbers are stored in $a_1, a_2, \cdots, a_n$;
then the following pseudocode gives one of the simplest implementations of
the insertion sort algorithm:

Below is the sample code in C/C++:

\begin{verbatim}
for (int i = 1; i <= n; i++)
    for (int j = i; j>=2; j--)
        if ( a[j] < a[j-1] ){
            int t = a[j-1];
            a[j-1] = a[j];
            a[j] = t;
        }
\end{verbatim}

Below is the sample code in Pascal:

\begin{verbatim}
for i:=1 to n do
    for j:=i downto 2 do
        if a[j]<a[j-1] then
            begin
                t:=a[i];
                a[i]:=a[j];
                a[j]:=t;
            end;
\end{verbatim}

To help Little Z understand insertion sort better, Little Z's teacher,
Teacher H, assigned the following homework:

Teacher H gives an array $a$ of length $n$, indexed from 1, in which all
elements are non-negative integers.
Little Z needs to support $Q$ operations on the array $a$.
There are two kinds of operations, with parameters as follows:

\texttt{1 x v}: this is the first kind of operation.
It changes the $x$-th element of $a$, that is, the value of $a_x$, to $v$.
It is guaranteed that $1 \le x \le n$ and $1 \le v \le 10^9$.
\textbf{Note that this operation changes the elements of the array;
the modified array is kept and affects subsequent operations.}

\texttt{2 x}: this is the second kind of operation.
Suppose Teacher H sorts the array $a$ according to
\textbf{the pseudocode above};
you need to tell Teacher H the position, in the new sorted array, of the
original $x$-th element of $a$, that is, $a_x$.
It is guaranteed that $1 \le x \le n$.
\textbf{Note that this operation does not change the elements of the
array; the sorted array is not kept and does not affect subsequent
operations.}

Teacher H does not like too many modifications, so he guarantees that the
number of operations of type 1 does not exceed 5000.

Little Z has never studied for programming competitions, so Little Z
cannot solve this problem.
He has come to you to help him solve it.

\subsection*{Input}

Read the input from the file \texttt{sort.in}.

The first line of input contains two positive integers $n, Q$,
denoting the length of the array and the number of operations.
It is guaranteed that $1 \le n \le 8\,000$ and $1 \le Q \le 2 \times 10^5$.

The second line of input contains $n$ space-separated non-negative
integers, where the $i$-th non-negative integer denotes $a_i$.
It is guaranteed that $1 \le a_i \le 10^9$.

The following $Q$ lines each contain $2 \sim 3$ positive integers
describing one operation;
the format of the operations is given in the description.

\subsection*{Output}

Write the output to the file \texttt{sort.out}.

For every query of type 2, output one line containing one positive integer,
the answer.

\subsection*{Sample 1 input}

\begin{verbatim}
3 4
3 2 1
2 3
1 3 2
2 2
2 3
\end{verbatim}

\subsection*{Sample 1 output}

\begin{verbatim}
1
1
2
\end{verbatim}

\subsection*{Explanation of sample 1}

Before the modification operation, suppose Teacher H performs one insertion
sort;
then the positions of the three elements of the original sequence after the
sort are $3, 2, 1$ respectively.

Before the modification operation, suppose Teacher H performs one insertion
sort;
then the positions of the three elements of the original sequence after the
sort are $3, 1, 2$ respectively.

Note that although $a_2 = a_3$ at this moment,
\textbf{we cannot regard them as the same element}.

\subsection*{Sample 2}

See \texttt{sort/sort2.in} and \texttt{sort/sort2.ans} in the contestant's
directory.

The data range of this test point is the same as that of test points
$1 \sim 2$.

\subsection*{Sample 3}

See \texttt{sort/sort3.in} and \texttt{sort/sort3.ans} in the contestant's
directory.

The data range of this test point is the same as that of test points
$3 \sim 7$.

\subsection*{Sample 4}

See \texttt{sort/sort4.in} and \texttt{sort/sort4.ans} in the contestant's
directory.

The data range of this test point is the same as that of test points
$12 \sim 14$.

\subsection*{Constraints}

For all test data,
$1 \le n \le 8\,000$, $1 \le Q \le 2 \times 10^5$, $1 \le x \le n$,
$1 \le v, a_i \le 10^9$.

For all test data, it is guaranteed that among all $Q$ operations, at most
5000 operations are of type one.

The additional restrictions and scores of the test points are shown in the
following table.

\begin{center}
\begin{tabular}{cccl}
\toprule
Test points & $n$ & $Q$ & Special property \\
\midrule
1,2,3,4     & $\le 10$    & $\le 10$            & None \\
\cmidrule(lr){1-3}
5,6,7,8,9   & $\le 300$   & $\le 300$           & \\
\cmidrule(lr){1-3}
10,11,12,13 & $\le 1\,500$ & $\le 1\,500$       & \\
\midrule
14,15,16    & $\le 8\,000$ & $\le 8\,000$       &
  All input $a_i, v$ are guaranteed to be pairwise distinct \\
\cmidrule(lr){1-1}\cmidrule(lr){4-4}
17,18,19    &             &                     & None \\
\cmidrule(lr){1-1}\cmidrule(lr){3-4}
20,21,22    &             & $\le 2 \times 10^5$ &
  All input $a_i, v$ are guaranteed to be pairwise distinct \\
\cmidrule(lr){1-1}\cmidrule(lr){4-4}
23,24,25    &             &                     & None \\
\bottomrule
\end{tabular}
\end{center}

In the original table the special property ``None'' spans test points
1--13,
the cell $n \le 8\,000$ spans test points 14--25,
the cell $Q \le 8\,000$ spans test points 14--19,
and the cell $Q \le 2 \times 10^5$ spans test points 20--25.

\section*{Problem 3: Network Connection
  (\begin{CJK}{UTF8}{gbsn}网络连接\end{CJK}, network)}

Time limit: 1 s.
Memory limit: 512 MB.
Total points: 100.

\subsection*{Description}

The TCP/IP protocol is an important protocol in the field of network
communication.
Your task today is to try to use this protocol to reproduce a simplified
network connection scenario.

In this problem, computers are divided into two major categories:
servers (Server) and clients (Client).
Servers are responsible for establishing connections, and clients are
responsible for joining connections.

There are $n$ computers in total that need to make network connections,
numbered $1 \sim n$.
These machines, in increasing order of their numbers, each issue one
operation of establishing a connection or joining a connection.

When trying to establish or join a connection, each machine must provide
an address string.
The address string provided by a server denotes the address at which it
tries to establish a connection;
the address string provided by a client denotes the address of the
connection it tries to join.

A well-formed address string must have the following features:

1. It must be of the form \texttt{a.b.c.d:e}, where $a, b, c, d, e$ are
all non-negative integers;

2. $0 \le a, b, c, d \le 255$ and $0 \le e \le 65535$;

3. none of $a, b, c, d, e$ may contain redundant leading zeros.

Correspondingly, an address string that is not well-formed may have the
following features:

1. It is not a string of the form \texttt{a.b.c.d:e},
for example it contains more than 3 characters \texttt{.}
or more than 1 character \texttt{:}, etc.;

2. one or more of the integers $a, b, c, d, e$ are outside the above
ranges;

3. one or more of the integers $a, b, c, d, e$ contain redundant leading
zeros.

For example, the address string \texttt{192.168.0.255:80} is well-formed,
but \texttt{192.168.0.999:80}, \texttt{192.168.00.1:10},
\texttt{192.168.0.1:088} and \texttt{192:168:0:1.233} are all not
well-formed.

If the address string provided by a server or a client when issuing its
operation is not well-formed, this operation is simply ignored.

In this problem, we assume that every address string satisfying the above
rules can take part in a normal connection;
you need not consider the actual meaning of each address string.

Due to network congestion and other reasons, two servers are not allowed to
use the same address string;
if this happens, the later server trying to establish a connection will
fail to establish a connection.
Apart from this, every server that provides a well-formed address string
successfully establishes a connection.

If a client providing a well-formed address tries to join a connection and
its address string is the same as the address string provided by some
earlier server that has successfully established a connection,
then this client successfully joins the connection, and we say that it is
connected to this server;
if no such server can be found, the client is considered to have failed to
join a connection.

Please note that although two different servers are not allowed to use the
same address string,
several clients using the same address string, and the same server being
connected to by several clients at the same time, are allowed.

Your task is simple:
given the type and address string of every computer, determine the
connection status of each computer.

\subsection*{Input}

Read the input from the file \texttt{network.in}.

Line 1 contains one positive integer $n$.

The following $n$ lines each contain 2 strings \texttt{op}, \texttt{ad},
giving the type and address string of each computer in increasing order of
their numbers.

Here \texttt{op} is guaranteed to be one of the strings \texttt{Server}
or \texttt{Client},
and \texttt{ad} is a non-empty string of length at most 25 consisting only
of digits, the character \texttt{.} and the character \texttt{:}.

The two strings on each line are separated by exactly one space, and there
are no extra spaces at the end of a line.

\subsection*{Output}

Write the output to the file \texttt{network.out}.

Output $n$ lines in total;
line $i$ contains one positive integer or one string denoting the
connection status of computer $i$.
Specifically:

If computer $i$ is a server, then:

1. if it provides a well-formed address string and successfully establishes
a connection, output the string \texttt{OK};

2. if it provides a well-formed address string but fails to establish a
connection because an earlier server has the same address string,
output the string \texttt{FAIL};

3. if the address string it provides is not well-formed, output the string
\texttt{ERR}.

If computer $i$ is a client, then:

1. if it provides a well-formed address string and successfully joins a
connection, output one positive integer: the number of the server this
client is connected to;

2. if it provides a well-formed address string but fails to join a
connection, output the string \texttt{FAIL};

3. if the address string it provides is not well-formed, output the string
\texttt{ERR}.

\subsection*{Sample 1 input}

\begin{verbatim}
5
Server 192.168.1.1:8080
Server 192.168.1.1:8080
Client 192.168.1.1:8080
Client 192.168.1.1:80
Client 192.168.1.1:99999
\end{verbatim}

\subsection*{Sample 1 output}

\begin{verbatim}
OK
FAIL
1
FAIL
ERR
\end{verbatim}

\subsection*{Explanation of sample 1}

Computer 1 is a server; it provides the well-formed address string
\texttt{192.168.1.1:8080} and successfully establishes a connection;

computer 2 is a server; it provides the same address string as computer 1
and fails to establish a connection;

computer 3 is a client; it provides the well-formed address string
\texttt{192.168.1.1:8080}, successfully joins a connection, and is
connected to server 1;

computer 4 is a client; it provides the well-formed address string
\texttt{192.168.1.1:80}, but no server can be found to connect to;

computer 5 is a client; the address string \texttt{192.168.1.1:99999} it
provides is not well-formed.

\subsection*{Sample 2 input}

\begin{verbatim}
10
Server 192.168.1.1:80
Client 192.168.1.1:80
Client 192.168.1.1:8080
Server 192.168.1.1:80
Server 192.168.1.1:8080
Server 192.168.1.999:0
Client 192.168.1.1.8080
Client 192.168.1.1:8080
Client 192.168.1.1:80
Client 192.168.1.999:0
\end{verbatim}

\subsection*{Sample 2 output}

\begin{verbatim}
OK
1
FAIL
FAIL
OK
ERR
ERR
5
1
ERR
\end{verbatim}

\subsection*{Sample 3}

See \texttt{network/network3.in} and \texttt{network/network3.ans} in the
contestant's directory.

\subsection*{Sample 4}

See \texttt{network/network4.in} and \texttt{network/network4.ans} in the
contestant's directory.

\subsection*{Constraints}

\begin{center}
\begin{tabular}{ccl}
\toprule
Test point numbers & $n \le$ & Special property \\
\midrule
1           & 10   & Properties 1, 2, 3 \\
\midrule
$2 \sim 3$  & 100  & Properties 1, 2, 3 \\
\cmidrule(lr){1-2}
$4 \sim 5$  & 1000 & \\
\cmidrule(lr){1-1}\cmidrule(lr){3-3}
$6 \sim 8$  &      & Properties 1, 2 \\
\cmidrule(lr){1-1}\cmidrule(lr){3-3}
$9 \sim 11$ &      & Property 1 \\
\cmidrule(lr){1-1}\cmidrule(lr){3-3}
$12 \sim 13$ &     & Property 2 \\
\cmidrule(lr){1-1}\cmidrule(lr){3-3}
$14 \sim 15$ &     & Property 4 \\
\cmidrule(lr){1-1}\cmidrule(lr){3-3}
$16 \sim 17$ &     & Property 5 \\
\cmidrule(lr){1-1}\cmidrule(lr){3-3}
$18 \sim 20$ &     & No special property \\
\bottomrule
\end{tabular}
\end{center}

In the original table the cell ``Properties 1, 2, 3'' spans test points
2--5,
and the cell $n \le 1000$ spans test points 4--20.

``Property 1'' is: all address strings are guaranteed to be well-formed;

``Property 2'' is: for any two different computers, if both are servers or
both are clients, then the address strings they provide are guaranteed to
be different;

``Property 3'' is: the number of every server is guaranteed to be smaller
than the numbers of all clients;

``Property 4'' is: all address strings are guaranteed to be of the form
\texttt{a.b.c.d:e}, where $a, b, c, d, e$ are all non-negative integers
not exceeding $10^9$ and without redundant leading zeros;

``Property 5'' is: all address strings are guaranteed to be of the form
\texttt{a.b.c.d:e}, where $a, b, c, d, e$ are all non-empty strings
consisting only of digits.

For $100\%$ of the data, it is guaranteed that $1 \le n \le 1000$.

\section*{Problem 4: Little Bear's Fruit Baskets
  (\begin{CJK}{UTF8}{gbsn}小熊的果篮\end{CJK}, fruit)}

Time limit: 1 s.
Memory limit: 512 MB.
Total points: 100.

\subsection*{Description}

In Little Bear's fruit shop there is a row of $n$ fruits.
Each fruit is either an apple or an orange, and the fruits are numbered
from left to right with the positive integers $1, 2, 3, \ldots, n$.
Fruits of the same kind that are arranged consecutively are called a
``block''.
Little Bear wants to pick this row of fruits into several fruit baskets,
in the following way:
each time, the leftmost fruit of every ``block'' is picked out
simultaneously, and together they form one fruit basket.
This operation is repeated until the fruits are used up.
Note that after each basket is picked, the ``blocks'' may change.
For example, after the only orange between two apple ``blocks'' is picked
away, the two apple ``blocks'' become one ``block''.
Please help Little Bear compute the fruits contained in each fruit basket.

\subsection*{Input}

Read the input from the file \texttt{fruit.in}.

The first line of input contains one positive integer $n$, denoting the
number of fruits.

The second line of input contains $n$ space-separated integers, where the
$i$-th number denotes the kind of the fruit numbered $i$:
1 stands for an apple and 0 stands for an orange.

\subsection*{Output}

Write the output to the file \texttt{fruit.out}.

Output several lines.

Line $i$ describes the fruit basket formed by the fruits picked out in the
$i$-th round.
Output the numbers of all fruits in this basket in increasing order,
with every two numbers separated by one space.

\subsection*{Sample 1 input}

\begin{verbatim}
12
1 1 0 0 1 1 1 0 1 1 0 0
\end{verbatim}

\subsection*{Sample 1 output}

\begin{verbatim}
1 3 5 8 9 11
2 4 6 12
7
10
\end{verbatim}

\subsection*{Explanation of sample 1}

This is the explanation of the first set of sample data.

At the beginning the fruits are 1 1 0 0 1 1 1 0 1 1 0 0, forming 6 blocks
in total.

In the first round of picking fruits to form a basket, the fruits numbered
1, 3, 5, 8, 9, 11 are picked out.

After that the remaining fruits are 1 0 1 1 1 0, forming 4 blocks in total.

In the second round of picking fruits to form a basket, the fruits numbered
2, 4, 6, 12 are picked out.

After that the remaining fruits are 1 1, forming only 1 block.

In the third round of picking fruits to form a basket, the fruit numbered 7
is picked out.

Finally the remaining fruit is 1, forming only 1 block.

In the fourth round of picking fruits to form a basket, the fruit numbered
10 is picked out.

\subsection*{Sample 2 input}

\begin{verbatim}
20
1 1 1 1 0 0 0 1 1 1 0 0 1 0 1 1 0 0 0 0
\end{verbatim}

\subsection*{Sample 2 output}

\begin{verbatim}
1 5 8 11 13 14 15 17
2 6 9 12 16 18
3 7 10 19
4 20
\end{verbatim}

\subsection*{Sample 3}

See \texttt{fruit/fruit3.in} and \texttt{fruit/fruit3.ans} in the
contestant's directory.

\subsection*{Constraints}

For $10\%$ of the data, $n \le 5$.

For $30\%$ of the data, $n \le 1000$.

For $70\%$ of the data, $n \le 50000$.

For $100\%$ of the data, $n \le 2 \times 10^5$.

\subsection*{Note}

Since the data size is large, C/C++ contestants are advised to use
\texttt{scanf} and \texttt{printf} for input and output.

\section*{Source}

CCF CSP-J 2021, second round (held 2021-10-23).
Problem pages on QOJ, in contest order:

\begin{itemize}
\item \texttt{https://qoj.ac/problem/18075} (candy)
\item \texttt{https://qoj.ac/problem/18076} (sort)
\item \texttt{https://qoj.ac/problem/18077} (network)
\item \texttt{https://qoj.ac/problem/18078} (fruit)
\end{itemize}

The statements were taken from the official PDF statement embedded on each
QOJ problem page;
the sample files were taken from the attachment archive of each page.

\end{document}
